\subsection*{Question 4.2 - Second-order model with a minimum-phase zero for the closed-loop pitch}

\paragraph{Step 1 - Model and cost function.}
We look for a second-order system with one minimum-phase zero (a zero in the left half-plane),
\begin{equation*}
G(s) = \frac{K\,(s + z)}{s^2 + 2\zeta\omega_n s + \omega_n^2}, \qquad z > 0 ,
\end{equation*}
with four parameters $p = [K,\ z,\ \omega_n,\ \zeta]$.
Only the amplitudes of the response were measured, so the fit is done on the gain (the magnitude of the frequency response),
\begin{equation*}
|G(j\omega; p)| = \frac{K\sqrt{\omega^2 + z^2}}{\sqrt{(\omega_n^2 - \omega^2)^2 + 4\zeta^2\omega_n^2\omega^2}} .
\end{equation*}
The measured gain at each tested frequency $\omega_i = 1, 5, 10, 15, 20$ rad/s is $g_i = A_{res,i}/A_{ref,i}$ (Table~\ref{tab:pitch-bode}).
The parameters are the ones that minimize the sum of the squared differences between the measured and the model gain,
\begin{equation*}
J(p) = \sum_{i=1}^{5} \big( g_i - |G(j\omega_i; p)| \big)^2 , \qquad p^\star = \arg\min_p J(p) ,
\end{equation*}
that is, written explicitly,
\begin{equation*}
J(K, z, \omega_n, \zeta) = \sum_{i=1}^{5} \left( \frac{A_{res,i}}{A_{ref,i}} - \frac{K\sqrt{\omega_i^2 + z^2}}{\sqrt{(\omega_n^2 - \omega_i^2)^2 + 4\zeta^2\omega_n^2\omega_i^2}} \right)^2 .
\end{equation*}
The magnitude does not depend on the sign of $z$, so the minimum-phase condition is imposed afterwards by taking $z > 0$.
For the same reason, the fit cannot distinguish the minimum-phase from the non-minimum-phase case by itself.
$J$ is nonlinear in $p$, so it is minimized numerically with the MATLAB function \texttt{nlinfit} (script \texttt{fit\_pitch\_bode.m}), as suggested in the handout.
By default, \texttt{nlinfit} solves this nonlinear least-squares problem with the Levenberg-Marquardt algorithm, which interpolates between Gauss-Newton and gradient descent, starting from an initial guess $p_0 = [1,\ 20,\ 7,\ 0.5]$.

\paragraph{Step 2 - Result.}
The result is
\begin{equation*}
K = 0.909, \qquad z = 21.2 \text{ rad/s}, \qquad \omega_n = 6.82 \text{ rad/s}, \qquad \zeta = 0.492 ,
\end{equation*}
that is,
\begin{equation*}
G(s) = \frac{0.909\, s + 19.3}{s^2 + 6.71\, s + 46.5} = \frac{0.909\,(s + 21.2)}{s^2 + 6.71\, s + 46.5} .
\end{equation*}
Other initial guesses (for example $[1, 5, 7, 0.5]$ and $[0.5, 50, 6, 1]$) lead to the same solution, so the fit does not depend on the starting point.

\paragraph{Step 3 - Quality of the fit.}
Table~\ref{tab:pitch-bode-fit} compares the measured and the fitted gains.
The root-mean-square error is $6.3 \times 10^{-3}$ in amplitude ratio.
The model follows the data at $1$, $5$ and $10$ rad/s (within $0.1$ dB), but at $15$ and $20$ rad/s the errors are $+0.8$ and $-1.2$ dB, which is $8$ to $15\%$ of the measured amplitude ratio.
These errors are larger than the uncertainty of the amplitudes (about $5\%$ at $20$ rad/s), and they have opposite signs: the model is below the data at $15$ rad/s and above it at $20$ rad/s, so the measured gain falls faster than the fitted model at high frequency.
The structure therefore does not fully capture the high-frequency roll-off.

\paragraph{Step 4 - Discussion.}
The identified parameters give the following picture of the closed-loop pitch.
\begin{itemize}
\item \emph{DC gain.} $G(0) = K z / \omega_n^2 = 19.26 / 46.54 \approx 0.414$ ($-7.7$ dB), consistent with the measured gain at $1$ rad/s ($0.418$).
It is below $1$ because the output is the angle in rad and the input is the raw command, so $0.414$ is the steady-state angle per unit of command.
\item \emph{Dynamics.} $\omega_n = 6.82$ rad/s and $\zeta = 0.492$, so the poles are complex, $s = -3.35 \pm 5.94\,j$.
The system is underdamped, with a small resonance peak of about $0.50$ ($-6.1$ dB) near $5$ rad/s, which agrees with the maximum of the measured data at $5$ rad/s.
The $-3$ dB bandwidth (relative to the DC gain) is about $9$ rad/s.
\item \emph{Zero.} The zero at $s = -z = -21.2$ rad/s lies well above $\omega_n$, so it only slows down the roll-off at high frequency, instead of shaping the response around the resonance.
It is therefore the least well determined parameter, because the data only reach $20$ rad/s.
\end{itemize}

\paragraph{Limitations.}
The model has four parameters and only five data points, so there is a single degree of freedom left, and the good fit at low frequency is not strong evidence that the structure is right.
The amplitude of the reference also changes from run to run, so any amplitude-dependent nonlinearity (saturation, dead zone, drone onboard controller) is mixed with the frequency dependence.
Finally, the Wi-Fi link and the onboard processing add a delay that is not in the model.
A delay does not change the magnitude, but it adds phase lag.
This can be checked against the step profile in Question 4.3.
